In this project we analyse a model for describing the trajectory of particles of pollutants in groundwater. A successful model for describing underground flows is given by the uncertain Darcy problem. Given a domain $D$ such that its boundary $\partial D$ is divided in three subsets $\Gamma_{\rm in}, \Gamma_{\rm out}, \Gamma_N$ such that 
\begin{equation*}
	\Gamma_{\rm in} \cup \Gamma_{\rm out} \cup \Gamma_N = \partial D, \quad \Gamma_{\rm in} \cap \Gamma_{\rm out} \cap \Gamma_N = \emptyset,
\end{equation*} 
the pressure and velocity fields $p$ and $u$ are given by the solution of the following Partial Differential Equation (PDE)
\begin{equation}
	\label{eq:IntroDarcy}
	\left \{
	\begin{aligned}
		u &= -A \nabla p, && \text{in } D, \\
		\nabla\cdot u &= f, && \text{in } D, \\
		p &= p_0, && \text{on } \Gamma_{\rm in},\\
		p &= 0, && \text{on } \Gamma_{\rm out}, \\
		\nabla p \cdot n &= 0, && \text{on } \Gamma_N,
	\end{aligned} \right.
\end{equation}
where $\Gamma_{\rm in}$, $\Gamma_{\rm out}$ are the inlet and outlet portions of the boundary of $D$, and an impermeability condition is imposed on $\Gamma_N$.
